\clearpage\thispagestyle{empty}
\pdfbookmark[0]{Sources for Instruments (Yearbook for 1822)}{instruments1822-sources}
\begin{center}\textbf{Sources for Instruments (Yearbook for 1822)}\end{center}
\noindent This AI translation follows \textit{Instrumente} in \href{https://gdz.sub.uni-goettingen.de/id/PPN236007467}{\textit{Werke}, vol.~VI} (Göttingen, 1874), pp.~424--428. In the \textit{Astronomisches Jahrbuch für das Jahr 1822}, edited by Johann Elert Bode and issued in Berlin in 1819, the \href{https://api.digitale-sammlungen.de/iiif/image/v2/bsb10538351_00241/full/1600,/0/default.jpg}{first-print heading} credits Gauss with the observations and instrument description and Dirksen with the comet’s orbital elements. The report occupies pp.~235--241 in the Bavarian State Library copy. Gauss writes in the first person and addresses the recipient directly.
\par\smallskip\noindent In the \href{https://api.digitale-sammlungen.de/iiif/image/v2/bsb10538351_00242/full/1600,/0/default.jpg}{first-print elements}, the time of perihelion is given on 27 June with a raised h. The translation follows the collected edition’s 28 June and raised n, retains the incomplete construction about the precision of the zenith distances, and gives both of Gauss’s statements about the value of \(a\). The \href{https://api.digitale-sammlungen.de/iiif/image/v2/bsb10538351_00246/full/1600,/0/default.jpg}{historical footnote on the two equations} is present in both editions. The collected edition’s reference for the weights points to the \hyperlink{native-415}{earlier report on Repsold’s meridian circle}. Translation and typesetting produced by AI.
\clearpage
